% Generated draft: requires mathematical and source review.
\documentclass[11pt]{article}
\usepackage[margin=25mm]{geometry}
\usepackage{fontspec}
\setmainfont{Latin Modern Roman}
\usepackage{amsmath,amssymb,mathtools,unicode-math}
\setmathfont{Latin Modern Math}
\usepackage{enumitem,needspace,xurl,hyperref}
\setlength{\parindent}{0pt}
\setlength{\parskip}{5pt}
\setlength{\emergencystretch}{3em}
\newcommand{\N}{\mathbb{N}}
\newcommand{\Z}{\mathbb{Z}}
\newcommand{\Q}{\mathbb{Q}}
\newcommand{\R}{\mathbb{R}}
\newcommand{\C}{\mathbb{C}}
\newcommand{\F}{\mathbb{F}}
\newcommand{\A}{\mathbb{A}}
\newcommand{\PP}{\mathbb{P}}
\newcommand{\E}{\mathbb{E}}
\newcommand{\eps}{\varepsilon}
\newcommand{\dd}{\,\mathrm{d}}
\begin{document}
\section*{UnsolvedMath ID: 30002008}
\textbf{Source problem code:} catalogue entry
\subsection{Definitions and mathematical statement}
Write \(\mathfrak g_n\), \((\mathfrak h^*/W)_n\), and \(\check G\) for the objects denoted by these symbols in the supplied source. Interpret \(Z(\mathfrak g_n)\) as the center of the universal enveloping algebra \(U(\mathfrak g_n)\). For a parameter \(\Lambda\in(\mathfrak h^*/W)_n\), let \(\chi_\Lambda\) denote the central character intended by the source's parameterization. Saying that the center acts through \(\Lambda\) means
\[
zv=\chi_\Lambda(z)v
\qquad\text{for every }z\in Z(\mathfrak g_n)\text{ and }v\in V.
\]
Let \(M_n(V)\) denote the induced critical-level module referred to in the source, and retain its notation
\[
Y_{n,-\Lambda}:=\operatorname{Op}^{\operatorname{ord}_n}_{\check G}(D)_{-\Lambda}
\]
for the proposed oper subscheme.

Scheme-theoretic support on a specified closed subscheme means that its defining ideal acts by zero through the relevant central action. In an affine presentation with ambient scheme \(X=\operatorname{Spec}A\), closed subscheme \(Y=\operatorname{Spec}(A/I)\), and an \(A\)-module \(M\), this means \(IM=0\), or equivalently that the action factors through \(A/I\).

\emph{Requested assertion, conditional on the missing source conventions.} For every index \(n\) in the intended range, every \(\Lambda\in(\mathfrak h^*/W)_n\), and every \(\mathfrak g_n\)-module \(V\) whose center acts through \(\chi_\Lambda\), the module \(M_n(V)\) is scheme-theoretically supported on \(Y_{n,-\Lambda}\). Thus, if \(I_{n,-\Lambda}\) is the defining ideal of this subscheme in an affine oper-space presentation carrying the specified critical central action, the assertion is
\[
I_{n,-\Lambda}\,M_n(V)=0.
\]
The sign \(-\Lambda\) is retained exactly as supplied.

\emph{Unresolved qualifications.} The record does not specify the base field, hypotheses on \(\mathfrak g\), the definition or indexing convention for \(\mathfrak g_n\), the range of \(n\), the meanings of \(\mathfrak h\), \(W\), and the subscript \(n\) on the quotient, or the central-character parameterization. It also omits the induction construction, the critical-level normalization, the relevant central action, the meanings of \(D\) and \(\operatorname{ord}_n\), and the scheme structure and parameter map defining the fiber labeled \(-\Lambda\). Consequently this is a conditional transcription of the intended assertion; a fully specified mathematical formulation requires these conventions.
\subsection{Short English statement}
The source asks whether prescribing the central character \(\Lambda\) of a \(\mathfrak g_n\)-module forces its induced critical-level module to be annihilated by the defining ideal of the oper subscheme labeled \(-\Lambda\). Essential constructions and conventions are missing. Supplied dated curation labels the problem partially solved and reports the case \(n=0\) as a theorem; these claims are not independently verified here.
\subsection{Sources}
\begin{itemize}
\item Supplied local UnsolvedMath registry record, permanent ID \(30002008\), title \emph{Oper Support of Induced Critical-Level Modules}, source problem code \texttt{catalogue entry}, in \emph{Oberwolfach Reports Open Problems}. The record gives proposed year 2012 and supplies no individual proposer.
\item \emph{Enveloping Algebras and Geometric Representation Theory} (2012), Oberwolfach Reports, report 13, pp.\ 733--809, DOI \texttt{10.4171/owr/2012/13}: \url{https://doi.org/10.4171/owr/2012/13}. The supplied background identifies Oberwolfach Reports as reports documenting mathematical workshops and their open-problem discussions. No precise internal location for this problem is supplied.
\item Supplied literature reference: \emph{Compatibility of the Feigin-Frenkel and Harish-Chandra Isomorphisms}, \url{https://arxiv.org/abs/1308.2995}.
\item Literature assessment dated \texttt{2026-08-21}, copied from supplied dated curation and not independently verified: the record labels the problem partially solved, reports the support statement as a theorem for \(n=0\), describes higher-jet center and Wakimoto constructions as supporting evidence, and states that the full scheme-theoretic oper-support assertion for arbitrary \(n\) is not established. The supplied background requests independent verification of time-sensitive status and cited claims.
\item The registry labels the statement \texttt{corrected\_verified} and records a manual semantic, typographic, and LaTeX review describing the cleanup as meaning-preserving, self-contained, and free of unsupported source-dependent additions. This is supplied review metadata; it does not resolve the missing mathematical definitions identified here.
\end{itemize}
\end{document}
